\section{模型架构：Qwen3-Next 与 Linear Attention}

\subsection{混合注意力架构}

Qwen3-Next 采用了 \textbf{3:1 Hybrid 架构}——三层 Linear Attention 配一层 Full Attention 交替堆叠。这一设计与月之暗面（Kimi）的架构``殊途同归''。

\begin{importantbox}{为什么要换架构？}
核心动机有三：
\begin{enumerate}
    \item \textbf{降低 KV Cache 成本}：Linear Attention 不需要维护完整的 KV Cache，大幅节省显存
    \item \textbf{支持超长上下文}：为 infinite long context 铺路
    \item \textbf{推理效率}：长序列场景下推理速度更快
\end{enumerate}
此外，Qwen 团队的 Attention Gate 机制获得了 NeurIPS Best Paper。
\end{importantbox}

\subsection{本章小结}

架构创新是支撑下一代能力（超长上下文、Agent 持续交互）的基础设施。Qwen3-Next 的 Hybrid Linear Attention 架构在保持模型质量的同时，为长程 Agent 交互提供了必要的效率保障。

